%! Representing a Categorical Variable — Unit 1, Lesson 1.2 — AP Practice
% EFFL: unscored AP-format practice. The graded practice is the DeltaMath assignment announced
% on the cover; this page and Check Your Understanding both carry NA in the cover's score column.
% Part (a) of the free response is the spiral item — it reaches back to Lesson 1.1 and makes
% students classify the variable before they are allowed to summarize it.
% Page lockstep with homework_key rests on \writelines{n} in the blank matching a terse
% \ansline{} in the key — keep key answers short enough not to wrap past the reserved lines.
\documentclass[10pt]{article}
\usepackage{apstats-article}
\usepackage{apstats-boxes}

\begin{document}

\pageheader{Unit 1, Lesson 1.2}{AP Practice}
% NAMESTRIP: no \namedateperiod here — the name row lives on the cover only.

\vspace{0.15in}

% Authored BYTE-IDENTICALLY in homework/ and homework_key/ — this box is what keeps the
% blank and the key the same number of pages.
\boxguard[8]
\begin{remindbox}
\small
This page is \textbf{not scored} --- your graded homework is in \textbf{DeltaMath}. Work this
if you are having trouble, and bring it to office hours or study hall for assistance.
\end{remindbox}

\vspace{0.12in}

\boxguard[26]
\begin{scenariobox}[The Recycling Audit]{navy}
\small
A city inspected recycling bins in two neighborhoods and recorded what each bin mainly
contained. Riverside had \textbf{250} bins inspected; Hillcrest had \textbf{500}.

\vspace{8pt}
\renewcommand{\arraystretch}{1.25}
\begin{tabularx}{\linewidth}{@{} l Y Y Y Y Y @{}}
\rowcolor{sky}
\textbf{Neighborhood} & \textbf{Paper} & \textbf{Plastic} & \textbf{Glass} &
\textbf{Contaminated} & \textbf{Total} \\
\hline
Riverside & 100 &  75 &  50 & 25 & 250 \\
Hillcrest & 190 & 160 & 110 & 40 & 500 \\
\hline
\end{tabularx}
\end{scenariobox}

\vspace{0.14in}

\boxguard[16]
\begin{headlinebox}{goldbg}
\small
\textbf{Section I --- Multiple Choice.} Circle the single best answer. Every answer choice is
about the context, not just the arithmetic --- read them all before choosing.
\end{headlinebox}

\vspace{0.10in}

\begin{enumerate}[leftmargin=1.9em, itemsep=0.13in]

  \item What percent of the Riverside bins inspected were contaminated?
        \begin{enumerate}[label=\textbf{(\Alph*)}, leftmargin=2.2em, itemsep=2pt]
          \item 2.5\%
          \item 8\%
          \item 10\%
          \item 25\%
          \item 40\%
        \end{enumerate}

  \item Which statement about contamination in the two neighborhoods is \textbf{correct}?
        \begin{enumerate}[label=\textbf{(\Alph*)}, leftmargin=2.2em, itemsep=2pt]
          \item Hillcrest is worse, because it had more contaminated bins.
          \item Riverside is worse, because it had more contaminated bins.
          \item Hillcrest had more contaminated bins, but Riverside had the higher
                contamination rate.
          \item Riverside had more contaminated bins, but Hillcrest had the higher
                contamination rate.
          \item The neighborhoods cannot be compared, because different numbers of bins
                were inspected.
        \end{enumerate}

  \item A relative frequency table is built for Riverside's four categories. Which must
        be true?
        \begin{enumerate}[label=\textbf{(\Alph*)}, leftmargin=2.2em, itemsep=2pt]
          \item The four entries add to 250.
          \item The four entries add to 1, or equivalently 100\%.
          \item The four entries add to 750, the total bins inspected in the city.
          \item Each entry must be larger than the matching entry for Hillcrest.
          \item The entries cannot be compared to Hillcrest's under any circumstances.
        \end{enumerate}

  \item The city wants one graph showing whether the two neighborhoods sort their recycling
        differently. Which display best supports that comparison?
        \begin{enumerate}[label=\textbf{(\Alph*)}, leftmargin=2.2em, itemsep=2pt]
          \item Side-by-side bars of the raw counts, because counts are the actual data.
          \item Side-by-side bars of the relative frequencies, because the neighborhoods had
                different numbers of bins.
          \item A single bar for each neighborhood's total.
          \item Bars of the raw counts with Hillcrest's halved to match Riverside's total.
          \item No graph can compare two neighborhoods.
        \end{enumerate}

\end{enumerate}

\vspace{0.14in}

\boxguard[16]
\begin{headlinebox}{goldbg}
\small
\textbf{Section II --- Free Response.} Show your reasoning. Every answer must be written
\emph{in the context of the problem} --- that is what earns credit on the AP exam.
\end{headlinebox}

\vspace{0.10in}

\begin{enumerate}[leftmargin=1.9em, itemsep=0.16in, resume]

  \item A school surveyed its 9th graders and its 12th graders about how they would prefer to
        spend an activity period. Among \textbf{300} ninth graders: Sports 150, Clubs 90,
        Study hall 60. Among \textbf{120} twelfth graders: Sports 42, Clubs 30, Study hall 48.

        \vspace{6pt}
        \textbf{(a)} Identify the individuals and the variable, and classify the variable as
        categorical or quantitative.\\[4pt]
        \writelines{2}

        \vspace{6pt}
        \textbf{(b)} Build the relative frequency table for the twelfth graders. Show each
        calculation.\\[4pt]
        \writelines{3}

        \vspace{6pt}
        \textbf{(c)} A student newspaper writes: \textbf{``Ninth graders like sports far more
        than seniors do --- 150 of them chose it, versus only 42 seniors.''} Explain what is
        wrong with this comparison, and give the corrected version.\\[4pt]
        \writelines{3}

        \vspace{6pt}
        \textbf{(d)} Identify one preference where the two grades differ most, and describe
        that difference in context using relative frequencies.\\[4pt]
        \writelines{2}

\end{enumerate}

\vspace{0.12in}

\boxguard[12]
\begin{extensionbox}
\small
Find a news story that compares two groups using raw counts --- two cities, two countries, two
schools. Write down the counts, then look up the sizes of the two groups and redo the
comparison as rates. Say whether the story's point survives.\\[4pt]
\writelines{2}
\end{extensionbox}

\vspace{0.12in}

\boxguard[10]
\begin{spiralbox}
\small
\textbf{Next lesson: Representing a Quantitative Variable with Graphs.} Counting worked here
because every value was one of a few category names. Next time the values are numbers that
barely repeat --- you cannot count how many students scored exactly 73.4 --- so we need a
different set of displays, and the one you choose changes what you see.
\end{spiralbox}

\end{document}
